% Part: lambda-calculus
% Chapter: introduction
% Section: church-rosser

\documentclass[../../../include/open-logic-section]{subfiles}

\begin{document}

\olfileid{lam}{int}{cr}
\olsection{चर्च--रॉसर गुणधर्म}

\begin{thm}
\ollabel{thm:church-rosser}
$M$, $N_1$ आणि $N_2$ ही पदे असून $M \red N_1$ आणि $M \red
N_2$ असे समजा. मग असे $P$ पद आहे की $N_1 \red P$ आणि $N_2 \red P$.
\end{thm}

\begin{cor}
$M$ चे न्यूनीकरण सामान्य रूपापर्यंत करता येते असे समजा. मग हे
सामान्य रूप एकमेव आहे.
\end{cor}

\begin{proof}
$M \red N_1$ आणि $M \red N_2$ असेल, तर मागील प्रमेयामुळे
असे पद~$P$ आहे की $N_1$ आणि $N_2$ या दोन्हींचे न्यूनीकरण
होऊन~$P$ मिळते. $N_1$ आणि~$N_2$ ही दोन्ही सामान्य रूपात
असतील, तर हे केवळ $N_1 \ident P \ident N_2$ असेल तेव्हाच शक्य आहे.
\end{proof}

शेवटी, दोन पदे $M$ आणि~$N$ यांचे न्यूनीकरण होऊन तेच पद मिळत
असेल, तर त्यांना \emph{$\beta$-सममूल्य}, किंवा फक्त
\emph{सममूल्य}, म्हणू; म्हणजेच असे $P$ असेल की $M
\red P$ आणि $N \red P$. हे $M \equal[\beta] N$ असे लिहितात.
\olref{thm:church-rosser} वापरून $\equal[\beta]$ हा
समतुल्यता संबंध असल्याचे तपासता येते; शिवाय, कोणत्याही $M$
आणि~$N$ साठी $M \red N$ किंवा $N \red M$ असेल, तर
$M \equal[\beta] N$. (खरे तर, हा गुणधर्म असलेला
$\equal[\beta]$ हा \emph{लहानात लहान} समतुल्यता संबंध आहे,
हे दाखवता येते.)

\end{document}
